% TV5 viết; TV4 cung cấp bằng chứng cho các nhận định.
\chapter{Đánh giá, hướng cải tiến và kết luận}
\label{ch:evaluation}
\owner{TV5 -- \MemberFive}{TV1--TV4}

\section{Hai điểm mạnh có bằng chứng}
\writingnote{Mỗi nhận định phải có điều kiện áp dụng và bằng chứng từ bài báo hoặc kết quả nhóm. Tách nguồn tác giả khỏi xác nhận thực nghiệm của nhóm.}
\subsection{Điểm mạnh thứ nhất}
\todo{Nhận định; điều kiện; bảng/hình/mục nguồn; ý nghĩa thực tế và giới hạn.}
\subsection{Điểm mạnh thứ hai}
\todo{Nhận định độc lập với điểm thứ nhất; kèm bằng chứng và phạm vi.}

\section{Hai hạn chế có thể kiểm chứng}
\subsection{Hạn chế thứ nhất}
\todo{Mô tả hạn chế, phép thử và bằng chứng nếu đã đo. Gợi ý trong plan: độ linh hoạt spline có thể ảnh hưởng train--test gap và chi phí; không kết luận overfitting chỉ vì G lớn.}
\subsection{Hạn chế thứ hai}
\todo{Mô tả hạn chế khác và phép thử. Gợi ý: ít tham số chưa chắc chạy nhanh; đối chiếu RMSE--thời gian, cùng máy và ngân sách optimizer.}
\begin{table}[htbp]\centering\small
\caption{Khung truy vết nhận định đánh giá.}\label{tab:claims}
\begin{tabularx}{\linewidth}{@{}YYYY@{}}\toprule
\textbf{Nhận định} & \textbf{Bằng chứng / nguồn} & \textbf{Đã kiểm?} & \textbf{Giới hạn}\\\midrule
Điểm mạnh 1 & \blank & \blank & \blank\\
Điểm mạnh 2 & \blank & \blank & \blank\\
Hạn chế 1 & \blank & \blank & \blank\\
Hạn chế 2 & \blank & \blank & \blank\\\bottomrule
\end{tabularx}\end{table}

\section{Hướng cải tiến gắn với hạn chế}
\todo{Nêu ít nhất một cải tiến cụ thể, cơ chế kỳ vọng, chi phí, đối chứng và tiêu chí thành công. Ví dụ định hướng trong plan: chọn G bằng validation hoặc dừng tăng lưới khi validation không cải thiện.}
\writingnote{Đánh dấu rõ đề xuất chưa làm, đã triển khai hay đã kiểm chứng. Nếu nhận là cải thiện thực nghiệm, phải có baseline/đối chứng cùng điều kiện và run ID.}

\section{Nhận định riêng của nhóm}
\todo{Trả lời khi nào nên cân nhắc KAN, khi nào dữ liệu hiện tại chưa đủ kết luận; tổng hợp lý thuyết, chất lượng xấp xỉ và chi phí thực nghiệm. Không biến chặn xấp xỉ thành bảo đảm huấn luyện hoặc tổng quát hóa.}

\section{Kết luận}
\todo{Tóm tắt điều đã giải thích, phạm vi tái lập thực tế, 2--3 kết quả định lượng có nguồn và hạn chế quan trọng nhất. Không đưa nhận định mới chưa có ở phần kết quả/đánh giá.}
\subsection{Công việc tiếp theo}
\todo{Liệt kê các bước mở rộng khả thi theo mức ưu tiên; phân biệt rõ với phần đã hoàn thành trong đồ án.}
